\section{Manual práctico de la sala de guardia}

\subsection{Estandarización de la documentación y el análisis retrospectivo de incidentes}

Otro énfasis de la semana 9: el trabajo del SRE no puede depender solo de la memoria; los incidentes, las hipótesis, los pasos de depuración y las razones de cada decisión deben quedar por escrito.

\begin{itemize}
    \item Cada alert, trace o metric spike que produzca un incident debe vincularse a un ticket
    \item Cada ticket registra el momento en que se disparó, el alcance del impacto y las acciones de mitigación que se tomaron después
    \item Al terminar el incidente se redacta un post-mortem que señale el root cause, las lessons learned y las future mitigations
    \item El AI SRE puede generar el borrador de manera automática, pero de todos modos requiere revisión humana, más contexto y verificación
\end{itemize}

\begin{warningbox}{No escribir el post-mortem equivale a no haber pasado realmente por el incident}
Una acción que solo hace que los números del monitoreo vuelvan a la normalidad puede apagar el fuego únicamente de manera temporal. La verdadera capacidad de resguardo proviene de convertir el incidente en conocimiento: la próxima vez que ocurra un problema parecido, el equipo puede saber de inmediato cómo se resolvió la vez anterior.
\end{warningbox}

\subsection{Las tres etapas del sistema de guardia}

Si se descompone la idea de la semana 9 en más detalle, se obtiene un modelo de madurez de tres etapas:

\begin{table}[H]
\centering
\begin{tabular}{p{0.22\textwidth} p{0.28\textwidth} p{0.28\textwidth}}
\toprule
Etapa & Tarea principal & Papel del AI SRE \\
\midrule
Etapa de observación & Unificar el monitoreo y los logs, definir umbrales de alerta & Ofrecer resúmenes en tiempo real y hacer highlight de patrones sospechosos \\
Etapa de respuesta & Ubicar con rapidez el root cause y ejecutar la mitigación & Sugerir comandos, citar el registro de operaciones y verificar la solución \\
Etapa de análisis retrospectivo & Redactar el post-mortem y actualizar el runbook & Organizar de manera automática la evidence matrix y recordar las follow-up tasks \\
\bottomrule
\end{tabular}
\caption{El papel del AI SRE en las tres etapas del sistema de guardia}
\end{table}

\subsection{Resumen de la sección}

Los incidentes no solo deben resolverse con rapidez, también deben convertirse en conocimiento escrito. La semana 9 nos enseña que, cuando el AI SRE se incorpora al flujo de guardia, su valor inicial proviene de observation + suggestion, y más adelante se extiende hacia controlled execution e incident learning.
